% Input-ready conventional-proof component.
% Requires amsmath, amssymb, amsthm and theorem/lemma/corollary environments.
% This component is independent of the version of the finite witness data.
% All labels in this file have the prefix fr:.
\providecommand{\per}{\operatorname{per}}

\section{An explicit integer positive-definite counterexample from a finite certificate}
\label{fr:section}

This section converts a finite positive-semidefinite endpoint certificate
with a specified margin into an integer positive-definite matrix and
two specified rational parameters in \((0,1)\). The estimates below are
conventional proofs; their explicit numerical regularization is separate
from the scope of the Lean certificates reported elsewhere.

For an ordered \(k\)-by-\(k\) positive-semidefinite matrix \(C\), write
\(D_1(C)=P'_1(C)\), \(c_k=\binom{k}{2}\), and
\[
 T(C)=\left\|\sum_{i<j}(v_i\wedge v_j)
                    \prod_{a\ne i,j}(v_a\cdot z)\right\|_{\mathrm F}^2,
 \qquad C=(v_i\cdot v_j)_{i,j}.
\]
Thus \(2D_1(C)=c_k\per C-T(C)\). The empty permanent is one.
For \(I\subseteq[n]\), let \(D_{\rm emb}(B[I])\) denote the derivative
obtained by allowing permutations on \(I\) and fixing every index in
\([n]\setminus I\), with inversions counted in the original order on
\([n]\).

\begin{lemma}[A uniform bound for embedded principal derivatives]
\label{fr:embedded-bound}
Let \(B\) be a real positive-semidefinite matrix of order \(n\ge2\),
let \(I\subseteq[n]\), and put \(k=|I|\), \(j=n-k\). Then
\begin{equation}\label{fr:embedded-estimate}
 |D_{\rm emb}(B[I])|
 \le \left\{\frac{c_k+c_k^2}{2}+\frac{jk^2}{2}\right\}\per B[I]
 \le n^4\per B[I].
\end{equation}
\end{lemma}

\begin{proof}
For \(C=B[I]\), define the permanental cofactor matrix
\[
 F_{ab}(C)=c_{ab}\per C(a\mid b),
\]
where \(C(a\mid b)\) deletes row \(a\) and column \(b\).
The matrix \((\per C(a\mid b))_{a,b}\) is a Gram matrix of the
products of all but one row form, in the Fischer inner product.
The Schur product theorem therefore shows that \(F(C)\) is positive
semidefinite. The singleton case of Lieb's permanent inequality
\cite{Lieb1966} gives
\(F_{aa}(C)\le\per C\), and hence
\begin{equation}\label{fr:F-bound}
 |F_{ab}(C)|\le\per C.
\end{equation}
For each fixed index \(u\notin I\), write
\[
 \operatorname{cut}_{u}F(C)
   =\sum_{\substack{a,b\in I\\a<u<b}}F_{ab}(C),
\]
with the labels on \(I\) inherited from \([n]\). Inserting the fixed
index \(u\) contributes twice the number of active indices mapped from
its left to its right. Summing over permutations gives the exact identity
\begin{equation}\label{fr:padding}
 D_{\rm emb}(C)=D_1(C)+2\sum_{u\notin I}\operatorname{cut}_{u}F(C).
\end{equation}
There are at most \(k^2/4\) pairs crossing each cut. Thus
\eqref{fr:F-bound} bounds the second term in \eqref{fr:padding} in
absolute value by \(jk^2\per C/2\).

Each of the \(c_k\) summands in the exterior polynomial defining
\(T(C)\) has squared Fischer norm
\[
 (c_{aa}c_{bb}-c_{ab}^2)\per C[\{a,b\}^c]
 \le c_{aa}c_{bb}\per C[\{a,b\}^c]\le\per C.
\]
The final inequality follows by applying the singleton form of Lieb's
inequality twice. The triangle inequality, or Cauchy--Schwarz, now gives
\(T(C)\le c_k^2\per C\). Therefore
\[
 |D_1(C)|\le\frac{c_k+c_k^2}{2}\per C.
\]
The same statement holds for \(k=0,1\), when both sides involving the
derivative vanish. These estimates prove the first inequality in
\eqref{fr:embedded-estimate}. For \(k\ge1\),
\((c_k+c_k^2)/2\le k^4/2\), while
\(jk^2/2\le n^3/2\le n^4/2\). This proves the second inequality;
the case \(k=0\) is immediate.
\end{proof}

\begin{theorem}[Explicit integer regularization]
\label{fr:integer-regularization}
Let \(n\ge6\), and let \(B\) be an integer real positive-semidefinite
matrix of order \(n\) whose diagonal entries are at least one. Suppose
that
\begin{equation}\label{fr:input-margin}
 T(B)>\frac{1001}{1000}\binom n2\per B.
\end{equation}
Set
\begin{equation}\label{fr:explicit-matrix}
 M=10^4n^3,\qquad A=MB+I_n,\qquad
 H_0=\max_i b_{ii},\qquad H=MH_0+1.
\end{equation}
Then \(A\) is an integer positive-definite matrix, its entries have
absolute value at most \(H\), and
\begin{equation}\label{fr:retained-margin}
 D_1(A)<-M^n\frac{n(n-5)}{20000}\per B<0.
\end{equation}
Its least eigenvalue is at least one, and is exactly one when \(B\)
is singular. If \(B\) has rank \(r<n\), then precisely \(n-r\)
eigenvalues of \(A\) equal one.
\end{theorem}

\begin{proof}
Put \(\eta=M^{-1}\). Expansion in the diagonal additions gives
\begin{equation}\label{fr:diagonal-expansion}
 D_1(B+\eta I_n)
 =D_1(B)+\sum_{\varnothing\ne J\subseteq[n]}
       \eta^{|J|}D_{\rm emb}(B[J^c]).
\end{equation}
The diagonal assumptions and the iterated singleton form of Lieb's
inequality imply
\(\per B[J^c]\le\per B\) for every \(J\). They also imply
\(\per B\ge\prod_i b_{ii}\ge1\). By
Lemma~\ref{fr:embedded-bound}, the sum in
\eqref{fr:diagonal-expansion} has absolute value at most
\[
 n^4\per B\bigl((1+\eta)^n-1\bigr).
\]
Since \(n\eta\le1/2\), the inequalities
\(\binom nj\le n^j\) and the geometric-series formula give
\[
 (1+\eta)^n-1
 \le\sum_{j=1}^\infty(n\eta)^j
 =\frac{n\eta}{1-n\eta}\le2n\eta.
\]
The perturbation in the derivative is consequently at most
\(2n^5\eta\per B=n^2\per B/5000\). On the other hand,
\eqref{fr:input-margin} and the endpoint identity imply
\[
 D_1(B)<-\frac{n(n-1)}{4000}\per B.
\]
Combining the last two estimates yields
\[
 D_1(B+\eta I_n)
 <-\left(\frac{n(n-1)}{4000}-\frac{n^2}{5000}\right)\per B
 =-\frac{n(n-5)}{20000}\per B<0.
\]
The derivative is homogeneous of degree \(n\) in the matrix entries,
so multiplication by \(M^n\) gives \eqref{fr:retained-margin}.
The remaining assertions follow from \(A=MB+I_n\), from the spectral
theorem, and from
\(|b_{ij}|\le\sqrt{b_{ii}b_{jj}}\le H_0\).
\end{proof}

A certificate may establish \eqref{fr:input-margin} by computing any
explicit lower bound \(U\le T(B)\) and checking the integer inequality
\(1000U>1001\binom n2\per B\). In particular, a sum of selected mutually
orthogonal parity-component norms is sufficient. The values of \(n\)
and \(H_0\) are read from the particular frozen row file; none of the
preceding estimates depends on a choice between witness versions.

\begin{corollary}[Specified rational points of strict decrease]
\label{fr:rational-interior}
Under the hypotheses of Theorem~\ref{fr:integer-regularization}, define
the positive integer
\begin{equation}\label{fr:Q}
 Q=4n!\,n^4H^n,
 \qquad q_1=1-\frac2Q,\qquad q_2=1-\frac1Q.
\end{equation}
Then \(0<q_1<q_2<1\), and
\begin{equation}\label{fr:strict-interior}
 P'_{q}(A)\le-\frac12\quad(q_1\le q\le1),
 \qquad P_{q_1}(A)>P_{q_2}(A).
\end{equation}
Thus the integer matrix \(A\) and the two rational parameters in
\eqref{fr:Q} form a fully specified positive-definite counterexample.
\end{corollary}

\begin{proof}
The polynomial \(P_q(A)\) has integer coefficients. Since
\(D_1(A)<0\), it follows that \(D_1(A)\le-1\).
For \(0\le q\le1\), termwise differentiation of the permutation
expansion and \(\operatorname{inv}(\sigma)\le\binom n2\le n^2\)
give
\[
 |P''_q(A)|\le n!\,n^4H^n=\frac Q4.
\]
Terms with fewer than two inversions contribute zero, so this bound
also holds at \(q=0\). For \(q\in[1-2/Q,1]\), the fundamental
theorem of calculus gives
\[
 P'_q(A)\le P'_1(A)+\frac Q4(1-q)\le-1+\frac12=-\frac12.
\]
Since \(Q>2\), the indicated interval lies in \((0,1]\), and
integration between \(q_1\) and \(q_2\) proves the strict inequality.
\end{proof}
